% Source: https://artofproblemsolving.com/community/c3h1684550_area_of_a_parallelogram
% Stable identity: aops-topic-1684550-q01
Let $AD=a$, so $AB=2a$ and $AM=MB=a$. Also $BC=AD=a$.

\begin{asy}
usepackage("amsmath");
size(6cm);
defaultpen(fontsize(10pt));
pair A=(0,0),B=(4,0),D=2*dir(62),C=B+D,M=(A+B)/2;
draw(A--B--C--D--cycle,linewidth(0.9));
draw(D--M--C,linewidth(0.8)+blue);
pair u=unit(D-M),v=unit(C-M);real r=0.28;
draw(M+r*u--M+r*(u+v)--M+r*v,linewidth(0.6));
label("$A$",A,SW);label("$B$",B,SE);label("$C$",C,NE);label("$D$",D,NW);label("$M$",M,S);
label("$a$",(A+M)/2,S);label("$a$",(M+B)/2,S);label("$a$",(A+D)/2,W);label("$a$",(B+C)/2,E);
label("\mbox{}",(min(currentpicture,user=true)+max(currentpicture,user=true))/2,black); // keeps the text after the figure black when LaTeX typesets the labels inline
\end{asy}

\textbf{$\angle DMC=90^\circ$.} Triangle $AMD$ is isosceles with $AM=AD$, so $\angle AMD=\angle ADM=\dfrac{180^\circ-\angle A}{2}$. Triangle $BMC$ is isosceles with $BM=BC$, so $\angle BMC=\dfrac{180^\circ-\angle B}{2}$. Since $\angle A+\angle B=180^\circ$ in a parallelogram,
\[
\angle AMD+\angle BMC=\dfrac{360^\circ-180^\circ}{2}=90^\circ,
\]
and therefore $\angle DMC=180^\circ-90^\circ=90^\circ$.

\textbf{Areas.} Triangle $DMC$ has base $DC$ and its vertex $M$ on the opposite side $AB$, so its area is half the area of the parallelogram. Since it is right-angled at $M$, its area is also $\dfrac{1}{2}MC\cdot MD$. Hence
\[
[ABCD]=2[DMC]=MC\cdot MD. \qquad\blacksquare
\]
